% !TEX root = ../dabirestan.tex
%% ============================================================
%%  فصل ۵ — پیشامدهای مستقل
%% ============================================================
\chapter{پیشامدهای مستقل}
\label{hs:ch5}

\chapterintro{\begin{itemize}
  \item بفهمید دو پیشامد «مستقل» یعنی چه و آن را آزمایش کنید؛
  \item «مستقل» را با «ناسازگار» اشتباه نگیرید؛
  \item احتمال آزمایش‌های تکراری (مثل چند پرتاب سکه یا چند پرتاب آزاد) را حساب کنید.
\end{itemize}}

%% ------------------------------------------------------------
\section{وقتی یک پیشامد چیزی دربارهٔ دیگری نمی‌گوید}
%% ------------------------------------------------------------

اگر یک سکه بیندازید و یک تاس، نتیجهٔ سکه هیچ تأثیری روی تاس ندارد. دانستن اینکه سکه شیر آمده، احتمال ۶ آمدن تاس را عوض نمی‌کند. چنین پیشامدهایی را \emph{مستقل} می‌نامیم.

\begin{definition}[label=hs:def:indep]{پیشامدهای مستقل}
دو پیشامد $A$ و $B$ را \term{مستقل}{independent} می‌نامیم اگر
\[ \P(A\cap B)=\P(A)\times\P(B). \]
اگر $\P(B)>0$، این همان است که بگوییم $\P(A\mid B)=\P(A)$؛ یعنی دانستن $B$ احتمال $A$ را عوض نمی‌کند. دو پیشامدی که مستقل نباشند \emph{وابسته} نام دارند.
\end{definition}

\begin{example}[label=hs:ex:coin-die]{سکه و تاس}
یک سکه و یک تاس می‌اندازیم. احتمال اینکه سکه شیر و تاس ۶ بیاید چقدر است؟
\begin{solution}
فضای نمونه $2\times6=12$ برآمد هم‌شانس دارد و فقط یکی از آن‌ها ($H$ و ۶) مطلوب است؛ پس احتمال $\tfrac1{12}$ است. می‌بینیم که $\tfrac1{12}=\tfrac12\times\tfrac16$؛ پس این دو پیشامد مستقل‌اند، همان‌طور که انتظار داشتیم.
\end{solution}
\end{example}

\begin{example}[label=hs:ex:die-indep]{آیا مستقل‌اند؟}
تاسی سالم پرتاب می‌شود. $A=\{2,4,6\}$ (زوج)، $B=\{1,2,3,4\}$ و $C=\{1,2,3\}$. آیا $A$ و $B$ مستقل‌اند؟ $A$ و $C$ چطور؟
\begin{solution}
$\P(A)=\tfrac12$، $\P(B)=\tfrac23$ و $A\cap B=\{2,4\}$ با احتمال $\tfrac13$. چون $\tfrac12\times\tfrac23=\tfrac13$، \textbf{$A$ و $B$ مستقل‌اند}: نیمی از اعضای $B$ زوج‌اند، همان‌طور که نیمی از کل اعداد زوج‌اند.\\
اما $A\cap C=\{2\}$ با احتمال $\tfrac16$ و $\P(A)\P(C)=\tfrac12\times\tfrac12=\tfrac14$. چون برابر نیستند، \textbf{$A$ و $C$ وابسته‌اند}: در $\{1,2,3\}$ فقط یک‌سوم اعداد زوج‌اند.
\end{solution}
\end{example}

\begin{example}[label=hs:ex:replace]{با برگرداندن و بدون برگرداندن}
جعبه‌ای ۵ مهرهٔ قرمز و ۳ مهرهٔ آبی دارد. دو مهره پشت سر هم برمی‌داریم. احتمال اینکه هر دو قرمز باشند را در دو حالت حساب کنید: (الف) مهرهٔ اول را \emph{برمی‌گردانیم}؛ (ب) \emph{برنمی‌گردانیم}.
\begin{solution}
(الف) با برگرداندن، جعبه دوباره مثل اول می‌شود و برداشت دوم مستقل از اولی است: $\tfrac58\times\tfrac58=\tfrac{25}{64}\approx0.39$.\\
(ب) بدون برگرداندن، برداشت دوم به اولی \emph{وابسته} است: $\tfrac58\times\tfrac47=\tfrac{5}{14}\approx0.36$.
\end{solution}
\end{example}

\begin{warning}
\textbf{مستقل با ناسازگار فرق دارد.} اگر $A$ و $B$ ناسازگار باشند (با هم رخ ندهند) و هر دو احتمال مثبت داشته باشند، آن‌ها \emph{وابسته‌اند}: اگر بدانیم $B$ رخ داده، مطمئن می‌شویم $A$ رخ نداده است! مثلاً در پرتاب تاس، «زوج آمدن» و «فرد آمدن» ناسازگارند ولی کاملاً به هم وابسته‌اند.
\end{warning}

%% ------------------------------------------------------------
\section{آزمایش‌های تکراری}
%% ------------------------------------------------------------

وقتی یک آزمایش را چند بار تکرار می‌کنیم و نتیجه‌ها روی هم اثر ندارند (مثل پرتاب‌های پیاپی سکه)، احتمال هر دنباله از نتیجه‌ها \emph{حاصل‌ضرب} احتمال‌هاست. این قانون برای بیش از دو پیشامد مستقل هم برقرار است:
\[ \P(A_1\cap A_2\cap\cdots\cap A_n)=\P(A_1)\times\P(A_2)\times\cdots\times\P(A_n). \]

\begin{example}[label=hs:ex:freethrow]{پرتاب آزاد}
بسکتبالیستی هر پرتاب آزاد را با احتمال $0.7$ گل می‌کند و پرتاب‌هایش مستقل از هم‌اند. اگر سه پرتاب انجام دهد، احتمال اینکه (الف) هر سه گل شود؛ (ب) هیچ‌کدام گل نشود؛ (ج) دست‌کم یکی گل شود؛ (د) دقیقاً دو تا گل شود چقدر است؟
\begin{solution}
(الف) $0.7\times0.7\times0.7=0.343$.\\
(ب) احتمال خطا در هر پرتاب $0.3$ است: $0.3^3=0.027$.\\
(ج) طبق قانون متمم: $1-0.027=0.973$.\\
(د) دنباله‌های با دقیقاً دو گل: «گل، گل، خطا»، «گل، خطا، گل» و «خطا، گل، گل». هر کدام احتمال $0.7\times0.7\times0.3=0.147$ دارد؛ پس جواب $3\times0.147=0.441$ است.
\end{solution}
\end{example}

\begin{example}[label=hs:ex:demere]{معمای شوالیه دومِره}
یک قمارباز فرانسوی در قرن هفدهم می‌پرسید: کدام محتمل‌تر است؟ (الف) در ۴ پرتاب یک تاس دست‌کم یک ۶ بیاید؛ (ب) در ۲۴ پرتاب دو تاس دست‌کم یک جفت‌شش بیاید.
\begin{solution}
(الف) احتمال «هیچ ۶ نیاید» در ۴ پرتاب مستقل برابر $(\tfrac56)^4$ است؛ پس احتمال مطلوب $1-(\tfrac56)^4\approx0.518$.\\
(ب) در هر پرتاب دو تاس، احتمال جفت‌شش $\tfrac1{36}$ است؛ پس احتمال مطلوب $1-(\tfrac{35}{36})^{24}\approx0.491$.\\
پس (الف) کمی محتمل‌تر است. پاسکال با همین محاسبه به این سؤال جواب داد؛ و این نامه‌نگاری‌ها آغاز علم احتمال شد!
\end{solution}
\end{example}

\begin{example}[label=hs:ex:lights]{چراغ‌های راهنمایی}
در مسیر مدرسه دو چراغ راهنمایی هست که مستقل از هم کار می‌کنند. احتمال سبز بودن چراغ اول $0.6$ و چراغ دوم $0.5$ است. احتمال اینکه (الف) هر دو سبز باشند؛ (ب) دست‌کم یکی قرمز باشد چقدر است؟
\begin{solution}
(الف) $0.6\times0.5=0.3$. \quad (ب) متمم «هر دو سبز» است: $1-0.3=0.7$.
\end{solution}
\end{example}

\begin{moreinfo}{سیستم‌های سری و موازی}
مهندسان برای بالا بردن اطمینان یک سیستم از قطعات «موازی» (یدکی) استفاده می‌کنند. فرض کنید هر قطعه مستقل از بقیه با احتمال $0.9$ سالم بماند.
\begin{itemize}
  \item \textbf{سه قطعهٔ سری} (همه باید سالم باشند، مثل حلقه‌های یک زنجیر): $0.9^3=0.729$.
  \item \textbf{سه قطعهٔ موازی} (یکی سالم بماند کافی است): $1-0.1^3=0.999$.
\end{itemize}
به همین دلیل هواپیماها چند موتور و چند سیستم کنترل موازی دارند.
\end{moreinfo}

\begin{warning}
یادآوری \textbf{خطای قمارباز}: چون پرتاب‌های سکه مستقل‌اند، احتمال شیر در پرتاب ششم همیشه $\tfrac12$ است؛ حتی اگر پنج پرتاب قبلی همه خط آمده باشند.
\end{warning}

%% ------------------------------------------------------------
\section*{خلاصهٔ فصل ۵}
\addcontentsline{toc}{section}{خلاصه و تمرین‌ها}
%% ------------------------------------------------------------

\begin{summary}[title={\leavevmode\boxicon{accent}{$\Sigma$}\ \ خلاصهٔ فصل}]
\begin{itemize}
  \item \textbf{مستقل:} $\P(A\cap B)=\P(A)\P(B)$، یا معادلاً $\P(A\mid B)=\P(A)$.
  \item ناسازگار $\ne$ مستقل؛ پیشامدهای ناسازگار (با احتمال مثبت) وابسته‌اند.
  \item برداشت \emph{با برگرداندن} مستقل است و \emph{بدون برگرداندن} وابسته.
  \item در آزمایش‌های تکراری مستقل، احتمال‌ها در هم ضرب می‌شوند؛ «دست‌کم یکی» $=1-$ «هیچ‌کدام».
\end{itemize}
\end{summary}

\subsection*{تمرین‌ها}

\begin{exercise}
سکه‌ای سالم ۴ بار پرتاب می‌شود. احتمال اینکه (الف) همه شیر بیاید؛ (ب) دست‌کم یک خط بیاید چقدر است؟
\answer{(الف) $\tfrac1{16}$ \quad (ب) $\tfrac{15}{16}$}
\end{exercise}

\begin{exercise}
دو تاس پرتاب می‌شوند. $A$: «تاس اول زوج است» و $B$: «مجموع دو تاس ۷ است». آیا $A$ و $B$ مستقل‌اند؟
\answer{بله؛ $\P(A\cap B)=\tfrac3{36}=\tfrac12\times\tfrac16$}
\end{exercise}

\begin{exercise}
اگر $A$ و $B$ مستقل باشند و $\P(A)=0.4$ و $\P(B)=0.5$، مقادیر $\P(A\cap B)$ و $\P(A\cup B)$ را بیابید.
\answer{$0.2$ و $0.7$}
\end{exercise}

\begin{exercise}
احتمال اینکه علی به یک سؤال تستی چهارگزینه‌ای به صورت شانسی درست جواب دهد $\tfrac14$ است. اگر به ۳ سؤال شانسی جواب دهد، احتمال اینکه (الف) همه را درست بزند؛ (ب) دست‌کم یکی را درست بزند چقدر است؟
\answer{(الف) $\tfrac1{64}$ \quad (ب) $1-\tfrac{27}{64}=\tfrac{37}{64}$}
\end{exercise}

\begin{exercise}
احتمال باران در هر یک از روزهای شنبه، یکشنبه و دوشنبه $0.2$ است و روزها مستقل فرض می‌شوند. احتمال اینکه در این سه روز دست‌کم یک روز باران بیاید چقدر است؟
\answer{$1-0.8^3=0.488$}
\end{exercise}
